% T23 — «Python»: блок кода с тёмно-синим фоном и подсветкой, светлая зона для условия.
% Профиль: информатика 11 класс, 5 задач (ЕГЭ задание 17/19).
\documentclass[a4paper,11pt]{article}

\usepackage{../../../../Lessons/_templates/preamble_worksheet}
\usepackage{../_shared/worksheet_check}

% --- Палитра: тёмно-синий код + индиго акцент ---
\definecolor{wcprimary}{HTML}{3730A3}
\definecolor{wcprimaryLight}{HTML}{C7D2FE}
\definecolor{wcprimaryBG}{HTML}{EEF2FF}
\definecolor{codebg}{HTML}{1E1B4B}
\definecolor{codefg}{HTML}{E0E7FF}
\definecolor{codekw}{HTML}{F472B6}      % розовый: ключевые слова
\definecolor{codenum}{HTML}{FBBF24}     % жёлтый: числа
\definecolor{codestr}{HTML}{34D399}     % зелёный: строки
\definecolor{codecmt}{HTML}{94A3B8}     % серый: комментарии
\definecolor{codeline}{HTML}{4338CA}    % акцент: номера строк

\renewcommand{\familydefault}{\sfdefault}

\geometry{a4paper, top=15mm, bottom=15mm, left=14mm, right=14mm}

% --- Хелперы для подсветки ---
\newcommand{\kw}[1]{\textcolor{codekw}{\bfseries #1}}
\newcommand{\num}[1]{\textcolor{codenum}{#1}}
\newcommand{\str}[1]{\textcolor{codestr}{#1}}
\newcommand{\cmt}[1]{\textcolor{codecmt}{\textit{#1}}}

% --- Заголовок: «IDE-стиль» с табами файлов ---
\renewcommand{\WorksheetTitle}[1]{%
  \par\noindent\begin{tikzpicture}
    \fill[codebg] (0,0) rectangle (\linewidth, -16mm);
    % три точки macOS
    \fill[red!70]    (4mm, -3mm) circle (1.2mm);
    \fill[yellow!80] (8mm, -3mm) circle (1.2mm);
    \fill[green!70]  (12mm, -3mm) circle (1.2mm);
    % таб
    \fill[wcprimary] (18mm, -1mm) rectangle (50mm, -6mm);
    \node[anchor=west, text=white, font=\ttfamily\footnotesize]
      at (19mm, -3.5mm) {worksheet.py};
    % заголовок
    \node[anchor=west, text=white, font=\sffamily\bfseries\fontsize{15pt}{18pt}\selectfont]
      at (4mm, -11mm) {#1};
    \node[anchor=east, text=codecmt, font=\ttfamily\footnotesize]
      at ([xshift=-4mm]\linewidth, -11mm) {Python 3.11 ·\ 5 заданий};
  \end{tikzpicture}\par\vspace{4mm}}

% --- Кодовый блок: тёмный фон, моноширинный текст ---
\newcommand{\CodeBlock}[1]{%
  \par\noindent\begin{tikzpicture}
    \node[fill=codebg, text=codefg, font=\ttfamily\footnotesize,
          inner xsep=4mm, inner ysep=3mm,
          text width=\dimexpr\linewidth-8mm\relax,
          rounded corners=3pt, anchor=north west]
      at (0,0) {#1};
  \end{tikzpicture}\par}

% --- TaskBox: светлая зона для условия + индиго ярлык ---
\renewcommand{\TaskBox}[2]{%
  \par\nopagebreak\vspace{2mm}\noindent\begin{tikzpicture}
    \node[fill=wcprimary, text=white, font=\sffamily\bfseries\footnotesize,
          rounded corners=2pt, inner xsep=4mm, inner ysep=1mm, anchor=west] (L) at (0,0)
      {def\ task\_#1():};
  \end{tikzpicture}\par\vspace{1mm}
  \noindent\begin{tikzpicture}
    \node[fill=wcprimaryBG, draw=wcprimaryLight, line width=0.4pt,
          rounded corners=4pt, inner xsep=4mm, inner ysep=3mm,
          text width=\dimexpr\linewidth-8mm\relax,
          anchor=north west] (B) at (0,0)
      {\small #2};
  \end{tikzpicture}\par\vspace{3mm}}

\SetAnswerFieldSize{34mm}{8mm}

\begin{document}

\WorksheetTitle{Алгоритмы и Python}
\par\noindent{\sffamily\small\color{wcprimary}\#\ ЕГЭ задание 17/19 \quad·\quad класс 11 \quad·\quad T23 \quad·\quad id=T23-python-2026-05-27}\par\vspace{3mm}
\FirstPageHeader

\TaskBox{1}{%
\textbf{Что выведет программа?}
\CodeBlock{%
\kw{def}\ f(n):\\
\ \ \ \ \kw{if}\ n\ <\ \num{2}:\\
\ \ \ \ \ \ \ \ \kw{return}\ \num{1}\\
\ \ \ \ \kw{return}\ n\ *\ f(n\ -\ \num{1})\\
\\
\kw{print}(f(\num{5}))\ \ \cmt{\#\ что выведет?}}\par
\vspace{2mm}\hfill\textbf{Ответ:}\ \AnswerField{1}}

\TaskBox{2}{%
\textbf{Сколько раз выполнится тело цикла?}
\CodeBlock{%
s = \num{0}\\
\kw{for}\ i\ \kw{in}\ \kw{range}(\num{1}, \num{50}):\\
\ \ \ \ \kw{if}\ i\ \%\ \num{3}\ ==\ \num{0}:\\
\ \ \ \ \ \ \ \ s\ +=\ \num{1}\\
\kw{print}(s)}\par
\vspace{2mm}\hfill\textbf{Ответ:}\ \AnswerField{2}}

\TaskBox{3}{%
\textbf{Найдите наименьшее $n$, при котором $f(n) > 1000$:}
\CodeBlock{%
\kw{def}\ f(n):\\
\ \ \ \ \kw{if}\ n\ ==\ \num{0}:\\
\ \ \ \ \ \ \ \ \kw{return}\ \num{1}\\
\ \ \ \ \kw{return}\ f(n\ -\ \num{1})\ *\ \num{2}\ +\ \num{1}}\par
\vspace{2mm}\hfill\textbf{Ответ:}\ \AnswerField{3}}

\TaskBox{4}{%
\textbf{Сколько различных значений примет переменная} \texttt{x}?
\CodeBlock{%
\kw{from}\ itertools\ \kw{import}\ product\\
result\ =\ \kw{set}()\\
\kw{for}\ a, b, c\ \kw{in}\ product([\num{1},\num{2},\num{3}], repeat=\num{3}):\\
\ \ \ \ x\ =\ a + b * c\\
\ \ \ \ result.add(x)\\
\kw{print}(\kw{len}(result))}\par
\vspace{2mm}\hfill\textbf{Ответ:}\ \AnswerField{4}}

\TaskBox{5}{%
\textbf{Что выведет программа?} (Рекурсивный обход)
\CodeBlock{%
\kw{def}\ g(n):\\
\ \ \ \ \kw{if}\ n\ <\ \num{1}: \kw{return}\ \str{""}\\
\ \ \ \ \kw{return}\ g(n // \num{2})\ +\ \kw{str}(n \% \num{2})\\
\\
\kw{print}(g(\num{45}))}\par
\vspace{2mm}\hfill\textbf{Ответ:}\ \AnswerField{5}}

\WorksheetQR{T23-python/L1/v1}

\end{document}
